\begin{abstract}
How far from finite can an amenable group be? Cornulier's sofic profile measures the least number of points on which a finite piece of a group
can be modelled by permutations at accuracy $1/r$. In earlier work we showed, by a counting criterion that can never give more
than $\exp(O(r))$, that Houghton's group has a finite piece with profile $\exp(\Omega(r^{0.129}))$; to our knowledge it was the first amenable
group shown to have superpolynomial profile. We construct a finitely presented elementary amenable group
with a finite piece whose profile is $\exp(r^{1+o(1)})$, with a lower bound certified by $25$ explicit relations: at least $\exp(c\,r/(\log r)^{12.75})$ and at
most $\exp(O(r\log r))$. The idea is to place copies of $S_3$ not on points but on configurations: Houghton's group moves the points of three rays,
finitely supported affine maps act on their binary configurations, and one copy of $S_3$ sits on each configuration. Then $N$ points carry $2^N$
copies, while any two copies are brought together by relations that touch only three points at a time. We reach this group through a ladder of
transport mechanisms, including an elementary amenable group at exponent $1/3$ and a subgroup of Thompson's group $V$ whose models, if they
exist, need $\exp(c\,r/\log^2r)$ points. Our group embeds in Brin's higher-dimensional Thompson group $3V$, so $3V$ has a chunk with finite
profile $\exp(r^{1+o(1)})$. Finally, the group becomes an explicit quantum channel on $\mathbb C^{873}$ whose repeated use is expensive.
A device that serves $n$ uses, releasing each output before the next input arrives, needs memory about $n/B$ when it may spend purity $B$, and
exact devices attain this trade-off up to polylogarithmic factors; with logarithmic purity it needs memory $n^{1-o(1)}$, and at the optimal
exchange rate its memory is $n^{1/2+o(1)}$, attained with purity of the same order. The engine is a one-round theorem that reads an approximate representation of the
group directly off the memory state of any such device. As a consequence, a local lattice process in $\nu$ dimensions whose bath starts
maximally mixed needs time about $\delta^{-1/\nu}$ to produce this channel to accuracy $\delta$. Autonomous dynamics can still produce such
channels: a bounded time-independent Hamiltonian, not local, coupled to a tracial bath, produces in the long run a channel with the same
memory-purity trade-off.
\end{abstract}
